\documentclass[11pt,a4paper]{article}

\newcommand{\doctitle}{Documento de Diseño Detallado}
\newcommand{\docsubtitle}{WARDEN - Wireless Attack Reproduction and Detection Environment}

\usepackage{warden-style}

% Paquete para bloques de codigo / pseudocodigo
\usepackage{listings}
\lstdefinestyle{wardencode}{
  basicstyle=\ttfamily\small,
  backgroundcolor=\color{wardenhighlight},
  frame=single,
  rulecolor=\color{wardenrule},
  framesep=4pt,
  framerule=0.4pt,
  xleftmargin=8pt,
  xrightmargin=8pt,
  breaklines=true,
  breakatwhitespace=false,
  keepspaces=true,
  showstringspaces=false,
  literate=
    {á}{{\'a}}1 {é}{{\'e}}1 {í}{{\'\i}}1 {ó}{{\'o}}1 {ú}{{\'u}}1
    {Á}{{\'A}}1 {É}{{\'E}}1 {Í}{{\'I}}1 {Ó}{{\'O}}1 {Ú}{{\'U}}1
    {ñ}{{\~n}}1 {Ñ}{{\~N}}1 {¿}{{?`}}1 {¡}{{!`}}1
}
\lstset{style=wardencode}

\begin{document}

% -- Portada -------------------------------------------------------------------
\wardencoverpage{0.1}{Abril 2026}{Borrador}

% -- Historial de Revisiones ---------------------------------------------------
\begin{wardenrevisions}
\revrow{0.1}{Abril 2026}{Santiago Valencia León}{Borrador inicial}
\revrow{0.1}{Abril 2026}{Daniel Yadir Perea Murillo}{Borrador inicial}
\revrow{0.1}{Abril 2026}{Sofia Valdez Londono}{Borrador inicial}
\end{wardenrevisions}

% -- Tabla de Contenidos -------------------------------------------------------
\tableofcontents
\newpage

% ==============================================================================
\section{Introducción}
% ==============================================================================

\subsection{Propósito}
Este documento describe el diseño detallado del sistema WARDEN, ofreciendo
suficiente información técnica para que la implementación pueda llevarse a
cabo de manera consistente y correcta por cualquier integrante del equipo.

A diferencia del Documento de Arquitectura, que describe la estructura
general del sistema y sus componentes a alto nivel, este documento entra
al detalle de cada componente: sus estructuras de datos, algoritmos,
interfaces y formatos de mensaje. Donde las decisiones de implementación
puedan refinarse durante la construcción, se indica explícitamente.

\subsection{Audiencia}
Este documento está dirigido principalmente a los integrantes del equipo
de desarrollo, que lo utilizarán como guía durante la implementación.
También sirve como evidencia para el profesor evaluador del nivel de
detalle del análisis técnico realizado antes de codificar.

\subsection{Alcance}
El documento cubre el diseño detallado de los dos subsistemas:

\begin{itemize}
  \item Subsistema Ofensivo (firmware ESP32): estructuras de datos,
    pseudocódigo de cada módulo, formato de frames generados, protocolo
    de consola serial, diseño del portal cautivo.
  \item Subsistema Defensivo (detector Python): estructuras de datos,
    algoritmos de análisis, formato de alertas, manejo de configuración.
\end{itemize}

Los detalles que no están cubiertos en este documento (porque dependen
de la implementación real) incluyen: nombres exactos de funciones,
manejo de excepciones específicas, optimizaciones internas. Estos
quedan a criterio del implementador, siempre que no contradigan el
diseño aquí descrito.

\subsection{Convención de Pseudocódigo}
A lo largo de este documento se utiliza pseudocódigo en español con
sintaxis simplificada. La intención es comunicar la lógica con claridad,
no producir código compilable. Cuando el pseudocódigo es trivial de
traducir a Python o C++, se usa una notación cercana al lenguaje destino.

\subsection{Referencias}
\begin{itemize}
  \item Documento WARDEN \emph{Documento de Arquitectura}.
  \item Documento WARDEN \emph{Especificación de Requerimientos de Software}.
  \item Documento WARDEN \emph{Documento de Casos de Uso}.
  \item Documento WARDEN \emph{Threat Model y Cadena de Ataque}.
  \item IEEE Std 802.11-2020.
\end{itemize}

% ==============================================================================
\section{Diseño del Subsistema Ofensivo}
% ==============================================================================

\subsection{Estructuras de Datos Principales}

\subsubsection{Configuración global}
La configuración del firmware se almacena en una estructura única que
todos los módulos consultan. Se inicializa con valores por defecto al
arranque y puede modificarse por consola serial.

\begin{lstlisting}
struct WardenConfig {
    uint8_t  bssid_objetivo[6];        // BSSID del router de laboratorio
    char     ssid_objetivo[33];        // SSID a clonar en Evil Twin
    uint8_t  canal;                    // Canal 802.11 (1-13)
    uint32_t duracion_beacon_flood;    // segundos, default 30
    uint32_t duracion_deauth;          // segundos, default 15
    uint32_t duracion_evil_twin;       // segundos, default 120
    uint32_t beacons_por_segundo;      // tasa de Beacon Flood, default 50
    char     prefijo_ssid_falso[16];   // para Beacon Flood, default "FakeNet"
    bool     bssid_validado;           // resultado del Validador Ético
};
\end{lstlisting}

\subsubsection{Estado de la cadena}
El Controlador de Cadena mantiene una máquina de estados explícita:

\begin{lstlisting}
enum EstadoCadena {
    IDLE,           // esperando comandos
    FASE_1_BEACON,  // ejecutando Beacon Flooding
    FASE_2_DEAUTH,  // ejecutando Deauthentication
    FASE_3_EVIL,    // ejecutando Evil Twin
    FINALIZADO,     // cadena completada, retornando a IDLE
    ERROR           // estado inalcanzable durante operacion normal
};

struct EstadoCadenaInfo {
    EstadoCadena    estado_actual;
    uint32_t        tiempo_inicio_fase;   // millis() al entrar a la fase
    uint32_t        beacons_emitidos;     // contador fase 1
    uint32_t        deauths_emitidos;     // contador fase 2
    uint32_t        clientes_conectados;  // contador fase 3
    uint32_t        credenciales_capturadas;  // contador fase 3
};
\end{lstlisting}

\subsubsection{Registro de credenciales capturadas}
Las credenciales recibidas en el portal cautivo se almacenan
exclusivamente en memoria volátil (RNF-SYS-005):

\begin{lstlisting}
struct CredencialCapturada {
    uint32_t   timestamp_millis;       // millis() al recibir
    char       cliente_ip[16];         // IP del cliente, formato ASCII
    char       usuario[64];            // texto recibido por POST
    char       password[64];           // texto recibido por POST
};

// Buffer circular de tamaño fijo, max 16 credenciales
CredencialCapturada credenciales[16];
int credenciales_count;
\end{lstlisting}

\subsection{Validador Ético: Algoritmo}

El Validador Ético es invocado cada vez que se intenta cambiar el BSSID
objetivo (UC-06). Implementa una validación de ``buena fe'' que rechaza
configuraciones obviamente externas al laboratorio.

\begin{lstlisting}
funcion validar_bssid(bssid):
    # Lista negra de prefijos OUI conocidos de fabricantes
    # comerciales mayoritarios (ejemplo: routers ISP comunes
    # en Colombia). Si el BSSID coincide con un OUI de esta
    # lista y NO coincide con la MAC del router de laboratorio
    # ya conocida, se rechaza.

    si bssid == bssid_router_laboratorio_conocido:
        retornar VALIDO

    si oui(bssid) en LISTA_NEGRA_OUI_ISP:
        retornar RECHAZADO_OUI_EXTERNO

    si bssid es broadcast (FF:FF:FF:FF:FF:FF):
        retornar RECHAZADO_BROADCAST

    si bssid es null (00:00:00:00:00:00):
        retornar RECHAZADO_INVALIDO

    # Para cualquier otro BSSID, se acepta solo si el operador
    # confirma explicitamente con el comando "set bssid X confirm"
    si comando_incluye_confirm:
        retornar VALIDO_CON_CONFIRMACION

    retornar REQUIERE_CONFIRMACION
\end{lstlisting}

\subsection{Módulo Beacon Flood: Pseudocódigo}

\begin{lstlisting}
funcion iniciar_beacon_flood(config):
    log_serial("Iniciando Beacon Flooding en canal " + config.canal)
    cambiar_modo_wifi_a_promiscuo()
    fijar_canal(config.canal)

    tiempo_fin = millis() + config.duracion_beacon_flood * 1000
    intervalo_us = 1_000_000 / config.beacons_por_segundo
    contador = 0

    mientras millis() < tiempo_fin y no hay_comando_stop():
        # Generar SSID y BSSID falsos
        ssid_falso = config.prefijo_ssid_falso + str(contador)
        bssid_falso = generar_bssid_aleatorio()

        # Construir frame beacon
        frame = construir_beacon(ssid_falso, bssid_falso, config.canal)

        # Emitir
        emitir_frame(frame)

        contador = contador + 1
        delay_microsegundos(intervalo_us)

    log_serial("Beacon Flooding finalizado, beacons emitidos: " + contador)
    retornar contador
\end{lstlisting}

\subsubsection{Estructura de un beacon falso}
Cada beacon emitido sigue el formato 802.11 estándar para frames de
gestión tipo \texttt{Beacon} (subtipo 0x08):

\begin{datatable}{L{4cm} L{2cm} L{9.3cm}}{Campo & Bytes & Valor}
\drow{Frame Control                & 2     & Tipo=Management (0), Subtipo=Beacon (8)}
\drow{Duration                     & 2     & 0x0000}
\drow{Destination Address          & 6     & FF:FF:FF:FF:FF:FF (broadcast)}
\drow{Source Address (BSSID)       & 6     & BSSID falso generado aleatoriamente}
\drow{BSSID                        & 6     & Mismo que Source Address}
\drow{Sequence Control             & 2     & Incremental}
\drow{Timestamp                    & 8     & millis() actual}
\drow{Beacon Interval              & 2     & 100 (TUs, equivale a 102.4 ms)}
\drow{Capability Information       & 2     & 0x0431 (red abierta, ESS)}
\drow{SSID Element                 & var   & Tipo=0, Length, SSID falso}
\drow{Supported Rates Element      & var   & Tipo=1, Length, tasas estándar}
\drow{DS Parameter Set             & 3     & Tipo=3, Length=1, Canal}
\end{datatable}

\subsection{Módulo Deauth: Pseudocódigo}

\begin{lstlisting}
funcion iniciar_deauth(config):
    log_serial("Iniciando Deauthentication contra BSSID " +
               formato_mac(config.bssid_objetivo))

    si NO config.bssid_validado:
        log_serial("ERROR: BSSID no validado, abortando")
        retornar ERROR_VALIDACION

    cambiar_modo_wifi_a_promiscuo()
    fijar_canal(config.canal)

    tiempo_fin = millis() + config.duracion_deauth * 1000
    contador = 0

    mientras millis() < tiempo_fin y no hay_comando_stop():
        # Construir frame deauth con BSSID legitimo como origen
        frame = construir_deauth(
            destino = FF:FF:FF:FF:FF:FF,           # broadcast a la red
            origen  = config.bssid_objetivo,        # suplantando AP
            bssid   = config.bssid_objetivo,
            razon   = 7                             # "Class 3 frame"
        )

        emitir_frame(frame)
        contador = contador + 1
        delay_milisegundos(50)  # ~20 deauths por segundo

    log_serial("Deauthentication finalizado, frames emitidos: " + contador)
    retornar contador
\end{lstlisting}

\subsubsection{Estructura de un frame deauth}

\begin{datatable}{L{4cm} L{2cm} L{9.3cm}}{Campo & Bytes & Valor}
\drow{Frame Control                & 2     & Tipo=Management (0), Subtipo=Deauthentication (12)}
\drow{Duration                     & 2     & 0x013A (314 us)}
\drow{Destination Address          & 6     & Cliente víctima específico o broadcast}
\drow{Source Address               & 6     & BSSID del router legítimo (suplantación)}
\drow{BSSID                        & 6     & BSSID del router legítimo}
\drow{Sequence Control             & 2     & Incremental}
\drow{Reason Code                  & 2     & 7 (clase 3 frame recibido en estado no asociado)}
\end{datatable}

\subsection{Módulo Evil Twin: Pseudocódigo}

\begin{lstlisting}
funcion iniciar_evil_twin(config):
    log_serial("Levantando Evil Twin con SSID " + config.ssid_objetivo)

    # 1. Iniciar AP modo open
    cambiar_modo_wifi_a_AP()
    iniciar_AP(
        ssid = config.ssid_objetivo,
        canal = config.canal,
        autenticacion = OPEN
    )

    # 2. Iniciar servidor DHCP
    iniciar_dhcp_server(
        ip_propia       = "10.0.0.1",
        rango_inicio    = "10.0.0.10",
        rango_fin       = "10.0.0.50",
        gateway         = "10.0.0.1",
        dns             = "10.0.0.1"
    )

    # 3. Iniciar servidor DNS local (resuelve todo a 10.0.0.1)
    iniciar_dns_server(ip_resolucion = "10.0.0.1")

    # 4. Iniciar servidor HTTP del portal cautivo
    iniciar_http_server(puerto = 80, handler = portal_handler)

    log_serial("Evil Twin activo. Esperando conexiones...")

    tiempo_fin = millis() + config.duracion_evil_twin * 1000
    mientras millis() < tiempo_fin y no hay_comando_stop():
        # El loop principal cede tiempo a las tareas de los servidores
        ceder_cpu()

    detener_http_server()
    detener_dns_server()
    detener_dhcp_server()
    detener_AP()

    log_serial("Evil Twin finalizado. Clientes: " +
               estado.clientes_conectados +
               ", credenciales: " + estado.credenciales_capturadas)
\end{lstlisting}

\subsubsection{Servidor DNS local: lógica}
El servidor DNS implementa una respuesta uniforme a cualquier consulta:

\begin{lstlisting}
funcion handler_dns(consulta):
    # consulta es un paquete DNS estandar
    dominio = extraer_dominio(consulta)

    log_serial("DNS query: " + dominio)

    # Construir respuesta apuntando a 10.0.0.1
    respuesta = construir_respuesta_dns(
        consulta = consulta,
        ip       = "10.0.0.1",
        ttl      = 60
    )

    enviar_respuesta(respuesta, cliente_origen)
\end{lstlisting}

\subsubsection{Portal cautivo: handler HTTP}

\begin{lstlisting}
funcion portal_handler(request):
    si request.metodo == "GET":
        # Servir la pagina HTML del portal
        servir_archivo("portal.html", content_type="text/html")
        retornar

    si request.metodo == "POST" y request.path == "/captura":
        # Procesar credenciales
        usuario = parsear_form(request.body, "usuario")
        password = parsear_form(request.body, "password")

        # Almacenar en buffer volatil
        credencial = CredencialCapturada {
            timestamp_millis = millis(),
            cliente_ip       = request.cliente_ip,
            usuario          = usuario,
            password         = password
        }
        credenciales.push(credencial)
        estado.credenciales_capturadas = estado.credenciales_capturadas + 1

        # Logar por serial
        log_serial("[CREDENCIALES CAPTURADAS]")
        log_serial("  IP cliente: " + request.cliente_ip)
        log_serial("  Usuario: " + usuario)
        log_serial("  Password: " + password)
        log_serial("  Timestamp: " + millis())

        # Responder con pagina de exito (demo)
        servir_string("<h1>Sesion iniciada (demo)</h1>",
                      content_type="text/html")
        retornar

    # Cualquier otra ruta o metodo, redirigir al portal
    redireccion(307, "/")
\end{lstlisting}

\subsubsection{Página HTML del portal cautivo}
El archivo \texttt{portal.html} embebido en el firmware es el siguiente
(simplificado):

\begin{lstlisting}
<!DOCTYPE html>
<html lang="es">
<head>
  <meta charset="UTF-8">
  <title>Re-autenticacion requerida</title>
</head>
<body style="font-family: sans-serif; max-width: 500px; margin: 50px auto;">
  <div style="border: 3px solid red; padding: 10px; margin-bottom: 20px;">
    <strong>WARDEN - DEMOSTRACION EDUCATIVA</strong><br>
    Este es un portal cautivo falso usado solo para fines educativos.<br>
    No ingreses credenciales reales.
  </div>
  <h1>Re-autenticacion requerida</h1>
  <p>Por favor inicia sesion para continuar.</p>
  <form action="/captura" method="POST">
    <p>Usuario: <input type="text" name="usuario" required></p>
    <p>Contrasena: <input type="password" name="password" required></p>
    <p><button type="submit">Iniciar sesion</button></p>
  </form>
</body>
</html>
\end{lstlisting}

\subsection{Diseño de la API REST de la ESP32}

El firmware ofensivo expone una API REST en \texttt{192.168.4.1:80}
accesible desde la red WiFi \texttt{WARDEN\_CONTROL}. Todos los endpoints
retornan JSON. Esta sección especifica los endpoints, sus parámetros y
sus esquemas de respuesta.

\subsubsection{AP de Control: configuración}

El AP \texttt{WARDEN\_CONTROL} se levanta al arranque de la ESP32
mientras no hay un ataque en curso:

\begin{lstlisting}
funcion iniciar_ap_control():
    config_ap = {
        ssid: "WARDEN_CONTROL",
        password: "warden-control-pwd",  # WPA2-PSK
        canal: 1,
        ip_propia: "192.168.4.1",
        rango_dhcp: "192.168.4.10 - 192.168.4.50"
    }
    iniciar_AP(config_ap)
    log_serial("AP de Control activo en 192.168.4.1")
\end{lstlisting}

\subsubsection{Endpoint: GET /status}

Retorna el estado actual del firmware. Pensado para que el panel
verifique la conectividad y lea contadores en vivo.

\begin{lstlisting}
GET /status

Respuesta 200 OK:
{
  "estado_cadena": "IDLE",   // IDLE | FASE_1 | FASE_2 | FASE_3 | FINALIZADO
  "fase_inicio_ms": 0,
  "uptime_ms": 543210,
  "contadores": {
    "beacons_emitidos": 0,
    "deauths_emitidos": 0,
    "clientes_evil_twin": 0,
    "credenciales_capturadas": 0
  },
  "ataque_activo": false
}
\end{lstlisting}

\subsubsection{Endpoint: GET /scan}

Ejecuta un escaneo activo de redes WiFi 2.4 GHz. Bloqueante durante
aproximadamente 10 segundos.

\begin{lstlisting}
GET /scan

Respuesta 200 OK:
{
  "duracion_seg": 10,
  "redes_encontradas": 8,
  "redes": [
    {
      "ssid": "LAB_WARDEN_UTP",
      "bssid": "e4:ab:89:d6:9b:80",
      "canal": 6,
      "rssi_dbm": -42,
      "cifrado": "WPA2"
    },
    {
      "ssid": "VECINO_WIFI",
      "bssid": "aa:bb:cc:dd:ee:ff",
      "canal": 11,
      "rssi_dbm": -78,
      "cifrado": "WPA2"
    },
    ...
  ]
}
\end{lstlisting}

\subsubsection{Endpoint: GET /clients}

Captura pasivamente frames dirigidos a un BSSID objetivo y retorna las
MACs detectadas. Bloqueante durante \texttt{duration\_seg} segundos
(default 30).

\begin{lstlisting}
GET /clients?bssid=e4:ab:89:d6:9b:80&duration=30

Respuesta 200 OK:
{
  "bssid_objetivo": "e4:ab:89:d6:9b:80",
  "duracion_captura_seg": 30,
  "clientes_detectados": 2,
  "clientes": [
    {
      "mac": "9c:ef:d5:f7:5a:e5",
      "frames_observados": 47,
      "primer_frame_ms": 1234,
      "ultimo_frame_ms": 28876
    },
    {
      "mac": "ac:5f:3e:11:22:33",
      "frames_observados": 12,
      "primer_frame_ms": 5678,
      "ultimo_frame_ms": 27432
    }
  ]
}

Respuesta 400 Bad Request:
{
  "error": "BSSID requerido y debe tener formato XX:XX:XX:XX:XX:XX"
}
\end{lstlisting}

\subsubsection{Endpoint: GET /oui-lookup}

Consulta la base de datos OUI embebida y retorna el fabricante asociado
a una MAC.

\begin{lstlisting}
GET /oui-lookup?mac=9c:ef:d5:f7:5a:e5

Respuesta 200 OK:
{
  "mac": "9c:ef:d5:f7:5a:e5",
  "oui": "9c:ef:d5",
  "fabricante": "Panda Wireless, Inc.",
  "encontrado": true
}

Respuesta 200 OK (cuando no se encuentra):
{
  "mac": "11:22:33:44:55:66",
  "oui": "11:22:33",
  "fabricante": "Fabricante desconocido",
  "encontrado": false
}
\end{lstlisting}

\subsubsection{Endpoint: GET /config y POST /config}

Permite consultar y modificar la configuración operativa del firmware.

\begin{lstlisting}
GET /config

Respuesta 200 OK:
{
  "bssid_objetivo": "e4:ab:89:d6:9b:80",
  "mac_victima": "9c:ef:d5:f7:5a:e5",
  "ssid_clonar": "LAB_WARDEN_UTP",
  "canal": 6,
  "duraciones_seg": {
    "beacon_flood": 30,
    "deauth": 15,
    "evil_twin": 120
  }
}

POST /config
Cuerpo de peticion:
{
  "bssid_objetivo": "e4:ab:89:d6:9b:80",
  "mac_victima": "9c:ef:d5:f7:5a:e5",
  "ssid_clonar": "LAB_WARDEN_UTP",
  "canal": 6,
  "duraciones_seg": {
    "beacon_flood": 30,
    "deauth": 15,
    "evil_twin": 120
  }
}

Respuesta 200 OK: configuracion aceptada
{
  "ok": true,
  "config": { ... echo de la config aplicada ... }
}

Respuesta 403 Forbidden: validador etico rechaza
{
  "ok": false,
  "error": "BSSID en lista negra de OUIs externos. Configuracion rechazada.",
  "codigo": "ETHICAL_VALIDATOR_REJECTED"
}
\end{lstlisting}

\subsubsection{Endpoint: POST /attack/start}

Inicia la cadena ofensiva con la configuración actual. Antes de
iniciar, vuelve a invocar al Validador Ético sobre la configuración
para garantizar la verificación.

\begin{lstlisting}
POST /attack/start
Cuerpo de peticion:
{
  "modo": "cadena_automatica"   // cadena_automatica | beacon | deauth | eviltwin
}

Respuesta 200 OK:
{
  "ok": true,
  "ataque_iniciado": "cadena_automatica",
  "fase_inicial": "FASE_1"
}

Respuesta 403 Forbidden:
{
  "ok": false,
  "error": "Configuracion no autorizada por validador etico"
}

Respuesta 409 Conflict:
{
  "ok": false,
  "error": "Hay un ataque en curso. Detenga primero con POST /attack/stop"
}
\end{lstlisting}

\subsubsection{Endpoint: POST /attack/stop}

\begin{lstlisting}
POST /attack/stop

Respuesta 200 OK:
{
  "ok": true,
  "fase_detenida": "FASE_2",
  "duracion_seg": 47.3
}

Respuesta 200 OK (cuando no hay ataque):
{
  "ok": true,
  "fase_detenida": null,
  "mensaje": "No habia ataque en curso"
}
\end{lstlisting}

\subsubsection{Endpoint: GET /credentials}

Retorna las credenciales ficticias capturadas durante la sesión actual
(buffer volátil en RAM).

\begin{lstlisting}
GET /credentials

Respuesta 200 OK:
{
  "total": 2,
  "credenciales": [
    {
      "timestamp_ms": 187432,
      "cliente_ip": "10.0.0.12",
      "usuario": "demo.user@warden.test",
      "password": "demo-password-12345"
    },
    {
      "timestamp_ms": 195203,
      "cliente_ip": "10.0.0.13",
      "usuario": "test@example.com",
      "password": "test123"
    }
  ]
}
\end{lstlisting}

\subsubsection{Endpoint: GET /events (Server-Sent Events)}

Stream de eventos en tiempo real para que el panel actualice su vista
sin polling agresivo. Implementación opcional con cabecera
\texttt{Content-Type: text/event-stream}.

\begin{lstlisting}
GET /events

Respuesta (stream continuo):
event: phase_change
data: {"fase": "FASE_2", "timestamp_ms": 30123}

event: client_connected
data: {"mac": "9c:ef:d5:f7:5a:e5", "ip": "10.0.0.12", "timestamp_ms": 145600}

event: credential_captured
data: {"usuario": "demo.user@warden.test", "ip": "10.0.0.12", "timestamp_ms": 145890}

event: phase_change
data: {"fase": "FINALIZADO", "timestamp_ms": 223456}
\end{lstlisting}

\subsection{Diseño del Frontend del Panel del Atacante}

El frontend está implementado como una SPA (Single Page Application)
ligera con HTML + JavaScript vanilla + Tailwind CSS via CDN.

\subsubsection{Estructura de archivos}

\begin{lstlisting}
attacker-panel/
+-- index.html              (entry point, layout con tabs)
+-- assets/
|   +-- panel.js            (logica principal del panel)
|   +-- api-client.js       (wrapper de fetch para la API REST)
|   +-- state.js            (estado global del panel)
|   +-- views/
|   |   +-- recon.js        (vista de reconocimiento)
|   |   +-- ethics.js       (vista de confirmacion etica)
|   |   +-- attack.js       (vista de ataque en vivo)
|   |   +-- summary.js      (vista de resumen)
|   +-- styles.css          (overrides minimos de Tailwind)
\end{lstlisting}

\subsubsection{Módulo api-client.js}

\begin{lstlisting}
// Cliente unificado de la API REST de la ESP32
const ESP32_BASE_URL = "http://192.168.4.1";

async function apiGet(path) {
  const res = await fetch(ESP32_BASE_URL + path);
  if (!res.ok) throw new Error("Error " + res.status);
  return await res.json();
}

async function apiPost(path, body) {
  const res = await fetch(ESP32_BASE_URL + path, {
    method: "POST",
    headers: { "Content-Type": "application/json" },
    body: JSON.stringify(body)
  });
  return { ok: res.ok, status: res.status, data: await res.json() };
}

// Wrappers semánticos
const api = {
  getStatus:       () => apiGet("/status"),
  scanNetworks:    () => apiGet("/scan"),
  detectClients:   (bssid) => apiGet("/clients?bssid=" + bssid),
  ouiLookup:       (mac) => apiGet("/oui-lookup?mac=" + mac),
  getConfig:       () => apiGet("/config"),
  setConfig:       (cfg) => apiPost("/config", cfg),
  startAttack:     (modo) => apiPost("/attack/start", { modo }),
  stopAttack:      () => apiPost("/attack/stop", {}),
  getCredentials:  () => apiGet("/credentials")
};
\end{lstlisting}

\subsubsection{Estructura del estado global (state.js)}

\begin{lstlisting}
// Estado global del panel, actualizado por las vistas
const state = {
  vista_actual: "recon",          // recon | ethics | attack | summary
  scan_resultados: [],
  red_seleccionada: null,         // objeto de la red elegida
  clientes_detectados: [],
  victima_seleccionada: null,     // MAC elegida
  etica_confirmada: false,
  ataque_estado: null,            // status del firmware
  credenciales: [],
  poll_interval: null
};

function setEstado(parcial) {
  Object.assign(state, parcial);
  renderizar(state.vista_actual);
}
\end{lstlisting}

\subsubsection{Vista en vivo del ataque (attack.js): polling}

Durante el ataque, el panel hace polling al endpoint \texttt{/status}
cada 1 segundo para refrescar contadores. Cuando detecta cambio de
fase a \texttt{FINALIZADO} o cuando el operador presiona detener,
transita a la vista de resumen.

\begin{lstlisting}
function iniciarMonitoreoAtaque() {
  state.poll_interval = setInterval(async () => {
    try {
      const status = await api.getStatus();
      setEstado({ ataque_estado: status });

      // Refrescar credenciales cada poll
      const creds = await api.getCredentials();
      setEstado({ credenciales: creds.credenciales });

      // Detectar fin del ataque
      if (status.estado_cadena === "FINALIZADO" || status.estado_cadena === "IDLE") {
        clearInterval(state.poll_interval);
        setEstado({ vista_actual: "summary" });
      }
    } catch (err) {
      // Conexion perdida (esperado durante Evil Twin)
      mostrarMensaje("Conexion temporal no disponible (fase Evil Twin activa)");
    }
  }, 1000);
}
\end{lstlisting}

\subsection{Protocolo de Consola Serial}

La interfaz serial de la ESP32 utiliza UART a 115200 baudios. Acepta los
siguientes comandos terminados en \texttt{\textbackslash n}:

\begin{datatable}{L{5cm} L{10.3cm}}{Comando & Acción}
\drow{\texttt{help}                              & Imprime la lista de comandos disponibles.}
\drow{\texttt{status}                            & Imprime el estado actual de la cadena, configuración y contadores.}
\drow{\texttt{set bssid XX:XX:XX:XX:XX:XX}       & Configura el BSSID objetivo. Pasa por el Validador Ético.}
\drow{\texttt{set bssid XX:XX:XX:XX:XX:XX confirm} & Variante con confirmación explícita para BSSIDs no en lista negra.}
\drow{\texttt{set ssid <NOMBRE>}                 & Configura el SSID a clonar para Evil Twin.}
\drow{\texttt{set channel N}                     & Configura el canal (1-13).}
\drow{\texttt{set duration <fase> <segundos>}    & Configura la duración de una fase (\texttt{beacon}, \texttt{deauth}, \texttt{eviltwin}).}
\drow{\texttt{start}                             & Inicia la cadena automática (UC-01).}
\drow{\texttt{beacon}                            & Ejecuta solo la fase de Beacon Flooding (UC-02).}
\drow{\texttt{deauth}                            & Ejecuta solo la fase de Deauthentication (UC-03).}
\drow{\texttt{eviltwin}                          & Ejecuta solo la fase de Evil Twin (UC-04).}
\drow{\texttt{stop}                              & Detiene cualquier ataque en curso (UC-05).}
\drow{\texttt{credentials}                       & Lista las credenciales capturadas en la sesión actual.}
\drow{\texttt{reset}                             & Limpia el buffer de credenciales y reinicia contadores.}
\end{datatable}

\subsubsection{Formato de mensajes de salida}
Todos los mensajes de salida del firmware siguen el formato:

\begin{lstlisting}
[NIVEL] [TIEMPO_MS] mensaje
\end{lstlisting}

Donde NIVEL es uno de: \texttt{INFO}, \texttt{WARN}, \texttt{ERROR},
\texttt{ALERT}, y TIEMPO\_MS es \texttt{millis()} desde el arranque.

Ejemplos:

\begin{lstlisting}
[INFO]  [12345] Iniciando Beacon Flooding en canal 6
[INFO]  [42345] Beacon Flooding finalizado, beacons emitidos: 1500
[ALERT] [89234] [CREDENCIALES CAPTURADAS]
[ALERT] [89234]   Usuario: demo.user@warden.test
[ALERT] [89234]   Password: demo-password-12345
[ERROR] [00100] BSSID no validado, comando rechazado
\end{lstlisting}

\subsection{Diseño de la API REST}

La ESP32 expone una API REST que el Panel del Atacante consume desde
\texttt{http://192.168.4.1}. Todos los endpoints retornan JSON. Los
errores se reportan con códigos HTTP estándar (400, 403, 404, 500) y
un cuerpo JSON con el formato \texttt{\{``error'': ``mensaje''\}}.

\subsubsection{Tabla de endpoints}

\begin{datatable}{L{4cm} L{1.5cm} L{9.8cm}}{Endpoint & Método & Descripción}
\drow{\texttt{/status}             & GET   & Estado actual del firmware (modo, fase, configuración mínima).}
\drow{\texttt{/config}             & GET   & Configuración completa actual.}
\drow{\texttt{/config}             & POST  & Actualiza la configuración. Pasa por Validador Ético.}
\drow{\texttt{/scan}               & GET   & Inicia escaneo activo de redes y retorna la lista.}
\drow{\texttt{/clients}            & GET   & Captura pasiva de clientes asociados a un BSSID. Acepta query \texttt{?bssid=XX}.}
\drow{\texttt{/oui-lookup}         & GET   & Resuelve un OUI a fabricante. Acepta query \texttt{?mac=XX}.}
\drow{\texttt{/attack/start}       & POST  & Inicia la cadena ofensiva con la configuración cargada.}
\drow{\texttt{/attack/stop}        & POST  & Detiene cualquier ataque en curso.}
\drow{\texttt{/attack/status}      & GET   & Estado actual del ataque (fase, contadores, tiempo restante).}
\drow{\texttt{/credentials}        & GET   & Lista de credenciales ficticias capturadas en la sesión.}
\drow{\texttt{/events}             & GET   & Server-Sent Events: stream de eventos en tiempo real.}
\end{datatable}

\subsubsection{Esquemas JSON}

\textbf{GET /status} — respuesta:

\begin{lstlisting}
{
  "uptime_ms": 458231,
  "modo": "control",
  "fase_actual": "idle",
  "bssid_objetivo": "e4:ab:89:d6:9b:80",
  "ssid_objetivo": "LAB_WARDEN_UTP",
  "canal": 6,
  "control_ap_activo": true,
  "api_disponible": true
}
\end{lstlisting}

El campo \texttt{modo} puede tomar los valores \texttt{control},
\texttt{ataque\_fase\_1}, \texttt{ataque\_fase\_2},
\texttt{ataque\_fase\_3}, \texttt{finalizado}.

\textbf{GET /scan} — respuesta:

\begin{lstlisting}
{
  "duracion_seg": 10,
  "redes_encontradas": 4,
  "redes": [
    {
      "ssid": "LAB_WARDEN_UTP",
      "bssid": "e4:ab:89:d6:9b:80",
      "canal": 6,
      "rssi": -42,
      "cifrado": "WPA2-PSK"
    },
    {
      "ssid": "Familia_Lopez",
      "bssid": "aa:bb:cc:11:22:33",
      "canal": 11,
      "rssi": -68,
      "cifrado": "WPA2-PSK"
    }
  ]
}
\end{lstlisting}

\textbf{GET /clients?bssid=e4:ab:89:d6:9b:80} — respuesta:

\begin{lstlisting}
{
  "bssid_objetivo": "e4:ab:89:d6:9b:80",
  "duracion_captura_seg": 30,
  "clientes_detectados": 2,
  "clientes": [
    {
      "mac": "9c:ef:d5:a1:b2:c3",
      "frames_observados": 47,
      "rssi_promedio": -55
    },
    {
      "mac": "88:e9:fe:01:02:03",
      "frames_observados": 12,
      "rssi_promedio": -71
    }
  ]
}
\end{lstlisting}

\textbf{GET /oui-lookup?mac=9c:ef:d5:a1:b2:c3} — respuesta:

\begin{lstlisting}
{
  "mac_consultada": "9c:ef:d5:a1:b2:c3",
  "oui": "9c:ef:d5",
  "fabricante": "Xiaomi Communications Co Ltd",
  "encontrado": true
}
\end{lstlisting}

Cuando el OUI no está en la base de datos:

\begin{lstlisting}
{
  "mac_consultada": "12:34:56:78:9a:bc",
  "oui": "12:34:56",
  "fabricante": null,
  "encontrado": false
}
\end{lstlisting}

\textbf{POST /config} — request:

\begin{lstlisting}
{
  "bssid_objetivo": "e4:ab:89:d6:9b:80",
  "mac_victima": "9c:ef:d5:a1:b2:c3",
  "ssid_a_clonar": "LAB_WARDEN_UTP",
  "canal": 6,
  "duracion_beacon_flood_seg": 30,
  "duracion_deauth_seg": 15,
  "duracion_evil_twin_seg": 120
}
\end{lstlisting}

Respuesta exitosa (HTTP 200):

\begin{lstlisting}
{
  "configuracion_aceptada": true,
  "validador_etico": "BSSID autorizado"
}
\end{lstlisting}

Respuesta con rechazo ético (HTTP 403):

\begin{lstlisting}
{
  "error": "BSSID rechazado por Validador Etico",
  "razon": "El OUI corresponde a un fabricante de routers ISP comerciales",
  "bssid_rechazado": "00:1a:2b:3c:4d:5e"
}
\end{lstlisting}

\textbf{POST /attack/start} — request:

\begin{lstlisting}
{
  "modo": "cadena_completa"
}
\end{lstlisting}

El campo \texttt{modo} acepta: \texttt{cadena\_completa},
\texttt{solo\_beacon}, \texttt{solo\_deauth}, \texttt{solo\_eviltwin}.

Respuesta exitosa (HTTP 200):

\begin{lstlisting}
{
  "ataque_iniciado": true,
  "modo": "cadena_completa",
  "tiempo_estimado_seg": 165
}
\end{lstlisting}

\textbf{GET /attack/status} — respuesta durante un ataque activo:

\begin{lstlisting}
{
  "ataque_activo": true,
  "fase_actual": "deauth",
  "tiempo_transcurrido_seg": 42,
  "tiempo_restante_fase_seg": 3,
  "contadores": {
    "beacons_emitidos": 1487,
    "deauths_emitidos": 234,
    "clientes_conectados_evil_twin": 0,
    "credenciales_capturadas": 0
  }
}
\end{lstlisting}

\textbf{GET /credentials} — respuesta:

\begin{lstlisting}
{
  "total": 1,
  "credenciales": [
    {
      "timestamp_ms": 89234,
      "cliente_ip": "10.0.0.15",
      "usuario": "demo.user@warden.test",
      "password": "demo-password-12345"
    }
  ]
}
\end{lstlisting}

\textbf{GET /events} — Server-Sent Events stream:

\begin{lstlisting}
event: fase_cambio
data: {"nueva_fase": "deauth", "timestamp_ms": 12345}

event: cliente_evil_twin_conectado
data: {"mac_cliente": "9c:ef:d5:a1:b2:c3", "ip_asignada": "10.0.0.15"}

event: credencial_capturada
data: {"usuario": "demo.user@warden.test", "ip": "10.0.0.15"}

event: ataque_finalizado
data: {"duracion_seg": 165, "credenciales_total": 1}
\end{lstlisting}

\subsubsection{Manejo de CORS y autenticación}

La API no implementa autenticación (la red \texttt{WARDEN\_CONTROL}
funciona como aislamiento de control de acceso). Sí incluye headers
\texttt{Access-Control-Allow-Origin: *} para permitir que el frontend
hospedado en cualquier origen pueda consumir la API sin restricciones
del navegador.

\subsection{Diseño del Panel del Atacante (Frontend)}

\subsubsection{Estructura de archivos}

\begin{lstlisting}
attacker-panel/
+-- index.html              (single page application)
+-- assets/
|   +-- panel.js            (logica principal de la app)
|   +-- api-client.js       (wrapper de fetch contra la ESP32)
|   +-- state.js            (gestion de estado global)
|   +-- styles.css          (overrides de Tailwind si aplica)
+-- README.md
\end{lstlisting}

\subsubsection{Vistas del Panel}

El panel es una SPA (single-page application) con tres vistas
principales gestionadas por un router simple en JavaScript vanilla:

\begin{datatable}{L{3cm} L{12.3cm}}{Vista & Contenido}
\drow{Reconocimiento  & Botones de ``Escanear redes'' y ``Detectar clientes''. Tablas de redes detectadas y clientes con sus fabricantes. Permite seleccionar red y víctima.}
\drow{Confirmación    & Resumen del ataque a lanzar. Casilla de confirmación ética obligatoria. Botón ``Lanzar ataque'' que solo se habilita tras marcar la casilla.}
\drow{Monitoreo       & Indicador de fase actual. Contadores en vivo. Lista de credenciales capturadas. Botón siempre visible de ``DETENER ATAQUE''. Resumen post-ataque cuando termina.}
\end{datatable}

\subsubsection{Cliente API: estructura}

\begin{lstlisting}
// api-client.js
const API_BASE = "http://192.168.4.1";

async function apiGet(path) {
    const response = await fetch(API_BASE + path);
    if (!response.ok) {
        const error = await response.json();
        throw new ApiError(response.status, error);
    }
    return await response.json();
}

async function apiPost(path, body) {
    const response = await fetch(API_BASE + path, {
        method: "POST",
        headers: {"Content-Type": "application/json"},
        body: JSON.stringify(body)
    });
    if (!response.ok) {
        const error = await response.json();
        throw new ApiError(response.status, error);
    }
    return await response.json();
}

// API funcional expuesta al resto del frontend
const api = {
    getStatus:      () => apiGet("/status"),
    scanNetworks:   () => apiGet("/scan"),
    detectClients:  (bssid) => apiGet(`/clients?bssid=${bssid}`),
    lookupOui:      (mac) => apiGet(`/oui-lookup?mac=${mac}`),
    setConfig:      (cfg) => apiPost("/config", cfg),
    startAttack:    (mode) => apiPost("/attack/start", {modo: mode}),
    stopAttack:     () => apiPost("/attack/stop", {}),
    getAttackStatus:() => apiGet("/attack/status"),
    getCredentials: () => apiGet("/credentials"),
    eventsStream:   () => new EventSource(API_BASE + "/events")
};
\end{lstlisting}

\subsubsection{Gestión de estado}

El frontend mantiene estado mínimo en memoria, sin localStorage (por la
restricción de no persistencia entre sesiones):

\begin{lstlisting}
// state.js
const state = {
    redObjetivo: null,        // {ssid, bssid, canal, rssi, cifrado}
    victimaSeleccionada: null,// {mac, fabricante}
    eticaConfirmada: false,
    ataqueActivo: false,
    faseActual: null,         // null | "beacon_flood" | "deauth" | "evil_twin"
    contadores: {
        beacons: 0, deauths: 0, clientes: 0, credenciales: 0
    },
    credencialesCapturadas: [],
    logsEvento: []            // ultimos 50 eventos del stream /events
};
\end{lstlisting}

\subsection{Mockup textual del Panel del Atacante}

A continuación se describe la estructura visual de la vista de
monitoreo (la más compleja) en formato de mockup textual:

\begin{lstlisting}
+--------------------------------------------------------------+
|  WARDEN - Panel del Atacante               [Estado: ATAQUE] |
+--------------------------------------------------------------+
|                                                              |
|  Objetivo:  LAB_WARDEN_UTP (e4:ab:89:d6:9b:80) Canal 6      |
|  Victima:   9c:ef:d5:a1:b2:c3 (Xiaomi Communications)        |
|                                                              |
|  +--------------------------------------------------------+  |
|  |  FASE 2 DE 3 - DEAUTHENTICATION                        |  |
|  |  [================================----] 80%           |  |
|  |  Tiempo restante: 3 segundos                           |  |
|  +--------------------------------------------------------+  |
|                                                              |
|  CONTADORES EN VIVO                                          |
|  - Beacons emitidos:        1487                             |
|  - Deauths emitidos:        234                              |
|  - Clientes en Evil Twin:   0                                |
|  - Credenciales capturadas: 0                                |
|                                                              |
|  CREDENCIALES CAPTURADAS                                     |
|  (Aun no hay credenciales)                                   |
|                                                              |
|  LOG DE EVENTOS (ultimos 5)                                  |
|  [12:34:01] Beacon Flooding iniciado                         |
|  [12:34:31] Beacon Flooding completado                       |
|  [12:34:31] Deauthentication iniciado                        |
|                                                              |
|              [  DETENER ATAQUE  ]  <-- siempre visible       |
|                                                              |
+--------------------------------------------------------------+
\end{lstlisting}

% ==============================================================================
\section{Diseño del Subsistema Defensivo}
% ==============================================================================

\subsection{Estructuras de Datos Principales}

\subsubsection{Configuración del detector}

\begin{lstlisting}
class DetectorConfig:
    iface: str                  # ej. "panda0", o None si modo PCAP
    pcap_file: str              # path al archivo, o None si modo live
    canal: int                  # canal monitoreado (1-13)
    bssid_protegido: bytes      # MAC del router de laboratorio (6 bytes)
    ssid_protegido: str         # SSID del router de laboratorio

    # Umbrales de detección (configurables)
    umbral_beacons_por_seg: int = 30        # tasa anómala
    ventana_beacon_seg: int = 5             # ventana de medición
    umbral_deauth_por_seg: int = 5          # tasa anómala
    ventana_deauth_seg: int = 3
    ventana_correlacion_seg: int = 60       # ventana para cadena

    bssid_lista_blanca: list[bytes]         # BSSIDs conocidos seguros
\end{lstlisting}

\subsubsection{Ventanas temporales}

Cada analizador mantiene una ventana móvil de eventos. Implementación
recomendada con \texttt{collections.deque}:

\begin{lstlisting}
from collections import deque

class VentanaTemporal:
    def __init__(self, duracion_segundos):
        self.duracion = duracion_segundos
        self.eventos = deque()  # [(timestamp, dato), ...]

    def agregar(self, timestamp, dato):
        self.eventos.append((timestamp, dato))
        self.purgar(timestamp)

    def purgar(self, ahora):
        # Elimina eventos mas viejos que la ventana
        while self.eventos and (ahora - self.eventos[0][0]) > self.duracion:
            self.eventos.popleft()

    def conteo(self):
        return len(self.eventos)

    def conteo_unicos(self, indice_dato):
        return len(set(e[1][indice_dato] for e in self.eventos))
\end{lstlisting}

\subsubsection{Registro de BSSIDs conocidos}
El Analizador Evil Twin mantiene un mapa SSID -> set de BSSIDs:

\begin{lstlisting}
class RegistroSsidBssid:
    def __init__(self):
        self.mapa = {}  # ssid (str) -> set[bssid (bytes)]

    def registrar(self, ssid, bssid):
        if ssid not in self.mapa:
            self.mapa[ssid] = set()
        self.mapa[ssid].add(bssid)

    def hay_duplicado(self, ssid, bssid):
        # Retorna True si el SSID ya existia con un BSSID distinto
        if ssid not in self.mapa:
            return False
        return bssid not in self.mapa[ssid] and len(self.mapa[ssid]) > 0
\end{lstlisting}

\subsection{Capturador de Frames: Diseño}

El Capturador abstrae la fuente de paquetes (vivo o PCAP) presentando
una interfaz uniforme:

\begin{lstlisting}
class Capturador:
    def __init__(self, config):
        self.config = config
        self.callback = None

    def registrar_callback(self, fn):
        # fn recibe cada frame como argumento
        self.callback = fn

    def iniciar(self):
        if self.config.iface:
            # Modo en vivo
            from scapy.all import sniff
            sniff(iface=self.config.iface,
                  prn=self.callback,
                  store=False,
                  monitor=True)
        elif self.config.pcap_file:
            # Modo offline
            from scapy.all import sniff
            sniff(offline=self.config.pcap_file, prn=self.callback)
\end{lstlisting}

\subsection{Despachador: Diseño}

\begin{lstlisting}
class Despachador:
    def __init__(self, config, analizadores):
        self.config = config
        self.analiz_beacon = analizadores['beacon']
        self.analiz_deauth = analizadores['deauth']
        self.analiz_eviltwin = analizadores['eviltwin']

    def procesar(self, frame):
        # Filtro de canal: descartar si el frame no es del canal monitoreado
        # (si el adaptador captura ruido de canales adyacentes)
        if not self.es_del_canal(frame, self.config.canal):
            return

        # Identificar tipo y enrutar
        if frame.haslayer(Dot11Beacon):
            self.analiz_beacon.procesar(frame)
            self.analiz_eviltwin.procesar(frame)
        elif frame.haslayer(Dot11Deauth):
            self.analiz_deauth.procesar(frame)
        # Otros tipos se ignoran
\end{lstlisting}

\subsection{Analizador de Beacon Flooding}

\textbf{Algoritmo:} mantiene una ventana de N segundos donde acumula
beacons únicos (por BSSID). Si la cantidad de BSSIDs únicos vistos en la
ventana supera el umbral configurado, emite alerta.

\begin{lstlisting}
class AnalizadorBeaconFlood:
    def __init__(self, config, reporter):
        self.config = config
        self.reporter = reporter
        self.ventana = VentanaTemporal(config.ventana_beacon_seg)
        self.alerta_emitida_recientemente = False
        self.ultimo_alerta = 0

    def procesar(self, frame):
        ahora = time.time()
        bssid = frame.addr3  # bytes
        ssid = frame.info.decode('utf-8', errors='replace')

        # Registrar este beacon en la ventana
        self.ventana.agregar(ahora, (bssid, ssid))

        # Contar BSSIDs unicos en ventana
        bssids_unicos = self.ventana.conteo_unicos(0)
        tasa = bssids_unicos / self.config.ventana_beacon_seg

        # Detectar si supera umbral, con cooldown anti-spam
        if tasa > self.config.umbral_beacons_por_seg:
            if ahora - self.ultimo_alerta > 5:  # cooldown 5s
                self.reporter.alerta(
                    severidad='ALERT',
                    tipo='BEACON_FLOOD',
                    mensaje=f"Tasa anomala de beacons: {tasa:.1f} BSSIDs unicos/s",
                    detalles={'bssids_unicos': bssids_unicos,
                              'ventana_seg': self.config.ventana_beacon_seg}
                )
                self.ultimo_alerta = ahora
\end{lstlisting}

\subsection{Analizador de Deauthentication}

\begin{lstlisting}
class AnalizadorDeauth:
    def __init__(self, config, reporter):
        self.config = config
        self.reporter = reporter
        self.ventana = VentanaTemporal(config.ventana_deauth_seg)
        self.ultimo_alerta = 0

    def procesar(self, frame):
        ahora = time.time()
        # Verificar si va dirigido al BSSID protegido
        bssid_origen = frame.addr2  # source address
        bssid_destino = frame.addr1  # destination
        bssid_field = frame.addr3   # BSSID en frame

        # Solo nos interesan deauths que involucran al BSSID protegido
        if bssid_field != self.config.bssid_protegido:
            return

        self.ventana.agregar(ahora, (bssid_origen, bssid_destino))

        conteo = self.ventana.conteo()
        tasa = conteo / self.config.ventana_deauth_seg

        if tasa > self.config.umbral_deauth_por_seg:
            if ahora - self.ultimo_alerta > 5:
                self.reporter.alerta(
                    severidad='ALERT',
                    tipo='DEAUTH',
                    mensaje=f"Deauths anomalos contra BSSID protegido: {tasa:.1f}/s",
                    detalles={'conteo': conteo,
                              'ventana_seg': self.config.ventana_deauth_seg,
                              'bssid_objetivo': formato_mac(bssid_field)}
                )
                self.ultimo_alerta = ahora
\end{lstlisting}

\subsection{Analizador de Evil Twin}

\begin{lstlisting}
class AnalizadorEvilTwin:
    def __init__(self, config, reporter):
        self.config = config
        self.reporter = reporter
        self.registro = RegistroSsidBssid()
        self.alertas_emitidas = set()  # (ssid, bssid_intruso)

    def procesar(self, frame):
        bssid = frame.addr3
        ssid = frame.info.decode('utf-8', errors='replace')

        # Si el SSID es el protegido pero el BSSID no es el protegido
        if ssid == self.config.ssid_protegido:
            if bssid != self.config.bssid_protegido:
                # Verificar lista blanca
                if bssid in self.config.bssid_lista_blanca:
                    return

                # Detectar Evil Twin
                clave_alerta = (ssid, bssid)
                if clave_alerta not in self.alertas_emitidas:
                    self.reporter.alerta(
                        severidad='CRITICAL',
                        tipo='EVIL_TWIN',
                        mensaje=f"BSSID {formato_mac(bssid)} anuncia SSID protegido",
                        detalles={
                            'ssid': ssid,
                            'bssid_intruso': formato_mac(bssid),
                            'bssid_legitimo': formato_mac(self.config.bssid_protegido)
                        }
                    )
                    self.alertas_emitidas.add(clave_alerta)

        # Siempre registrar SSID-BSSID para futura referencia
        self.registro.registrar(ssid, bssid)
\end{lstlisting}

\subsection{Correlador de Cadena}

El Correlador recibe las alertas de los tres analizadores y mantiene una
ventana temporal para detectar la cadena ofensiva completa.

\begin{lstlisting}
class CorreladorCadena:
    def __init__(self, config, reporter):
        self.config = config
        self.reporter = reporter
        self.alertas_recientes = {
            'BEACON_FLOOD': None,
            'DEAUTH': None,
            'EVIL_TWIN': None
        }
        self.cadena_emitida = False

    def registrar_alerta(self, tipo, timestamp):
        # Llamado por el Reporter al emitir cada alerta
        if tipo in self.alertas_recientes:
            self.alertas_recientes[tipo] = timestamp
            self.evaluar_cadena(timestamp)

    def evaluar_cadena(self, ahora):
        if self.cadena_emitida:
            return  # ya emitida, no repetir

        # Verificar si las tres alertas estan presentes y dentro de ventana
        timestamps = list(self.alertas_recientes.values())
        if any(t is None for t in timestamps):
            return

        ventana = self.config.ventana_correlacion_seg
        if (ahora - min(timestamps)) > ventana:
            return  # alguna alerta esta fuera de ventana

        # Verificar orden esperado: Beacon Flood -> Deauth -> Evil Twin
        t_beacon = self.alertas_recientes['BEACON_FLOOD']
        t_deauth = self.alertas_recientes['DEAUTH']
        t_eviltwin = self.alertas_recientes['EVIL_TWIN']

        if t_beacon <= t_deauth <= t_eviltwin:
            self.reporter.alerta(
                severidad='CRITICAL',
                tipo='CADENA_OFENSIVA',
                mensaje="Cadena ofensiva WARDEN detectada (las 3 fases en secuencia)",
                detalles={
                    't_beacon_flood': t_beacon,
                    't_deauth': t_deauth,
                    't_evil_twin': t_eviltwin,
                    'duracion_total_seg': t_eviltwin - t_beacon
                }
            )
            self.cadena_emitida = True
\end{lstlisting}

\subsection{Reporter: Formato de Alertas}

Todas las alertas siguen un formato uniforme y legible:

\begin{lstlisting}
[YYYY-MM-DDThh:mm:ss] [SEVERIDAD] TIPO_ATAQUE: mensaje
    detalle_1: valor_1
    detalle_2: valor_2
    ...
\end{lstlisting}

\textbf{Ejemplos reales:}

\begin{lstlisting}
[2026-04-25T14:23:01] [ALERT] BEACON_FLOOD: Tasa anomala de beacons: 87.4 BSSIDs unicos/s
    bssids_unicos: 437
    ventana_seg: 5

[2026-04-25T14:23:34] [ALERT] DEAUTH: Deauths anomalos contra BSSID protegido: 18.3/s
    conteo: 55
    ventana_seg: 3
    bssid_objetivo: e4:ab:89:d6:9b:80

[2026-04-25T14:24:02] [CRITICAL] EVIL_TWIN: BSSID aa:bb:cc:dd:ee:ff anuncia SSID protegido
    ssid: LAB_WARDEN_UTP
    bssid_intruso: aa:bb:cc:dd:ee:ff
    bssid_legitimo: e4:ab:89:d6:9b:80

[2026-04-25T14:24:03] [CRITICAL] CADENA_OFENSIVA: Cadena ofensiva WARDEN detectada (las 3 fases en secuencia)
    t_beacon_flood: 2026-04-25T14:23:01
    t_deauth: 2026-04-25T14:23:34
    t_evil_twin: 2026-04-25T14:24:02
    duracion_total_seg: 61.3
\end{lstlisting}

\subsection{Manejador de Sesión}

\begin{lstlisting}
class ManejadorSesion:
    def __init__(self, config):
        self.config = config
        self.tiempo_inicio = None
        self.contador_alertas = {
            'BEACON_FLOOD': 0, 'DEAUTH': 0,
            'EVIL_TWIN': 0, 'CADENA_OFENSIVA': 0
        }
        self.frames_procesados = 0
        self.capturador = None

    def iniciar(self):
        # Validaciones previas
        si self.config.iface:
            verificar_modo_monitor(self.config.iface)
        si self.config.pcap_file:
            verificar_archivo_existe(self.config.pcap_file)

        # Construir el pipeline
        reporter = Reporter(self)
        correlador = CorreladorCadena(self.config, reporter)
        reporter.set_correlador(correlador)

        analiz = {
            'beacon': AnalizadorBeaconFlood(self.config, reporter),
            'deauth': AnalizadorDeauth(self.config, reporter),
            'eviltwin': AnalizadorEvilTwin(self.config, reporter)
        }
        despachador = Despachador(self.config, analiz)
        self.capturador = Capturador(self.config)
        self.capturador.registrar_callback(self.on_frame_wrap(despachador))

        # Configurar manejador de SIGINT
        signal.signal(signal.SIGINT, self.shutdown)

        # Iniciar
        self.tiempo_inicio = time.time()
        log_info("WARDEN Detector iniciado.")
        self.capturador.iniciar()  # bloqueante hasta SIGINT o EOF

    def on_frame_wrap(self, despachador):
        def wrapper(frame):
            self.frames_procesados += 1
            despachador.procesar(frame)
        return wrapper

    def registrar_alerta_emitida(self, tipo):
        self.contador_alertas[tipo] += 1

    def shutdown(self, signum, frame):
        # Manejador de SIGINT
        duracion = time.time() - self.tiempo_inicio
        print("\n\n=== RESUMEN DE SESION WARDEN ===")
        print(f"Duracion: {duracion:.1f} segundos")
        print(f"Frames procesados: {self.frames_procesados}")
        print(f"Alertas emitidas:")
        for tipo, cuenta in self.contador_alertas.items():
            print(f"  {tipo}: {cuenta}")
        sys.exit(0)
\end{lstlisting}

\subsection{Argumentos de Línea de Comandos}

El detector se invoca con la siguiente sintaxis:

\begin{lstlisting}
python3 detector.py [opciones]

Opciones obligatorias (uno de los dos):
  --iface <nombre>      Captura en vivo desde la interfaz indicada
                        (debe estar en modo monitor)
  --pcap <archivo>      Analiza un archivo PCAP previamente capturado

Opciones de configuracion:
  --bssid <MAC>         BSSID del router de laboratorio (obligatorio)
  --ssid <NOMBRE>       SSID del router de laboratorio (obligatorio)
  --channel <N>         Canal monitoreado (1-13, default 6)

Opciones de umbrales (opcionales):
  --umbral-beacons <N>  BSSIDs unicos por segundo para alertar (default 30)
  --ventana-beacon <S>  Ventana en segundos (default 5)
  --umbral-deauth <N>   Deauths por segundo para alertar (default 5)
  --ventana-deauth <S>  Ventana en segundos (default 3)
  --ventana-corr <S>    Ventana de correlacion para cadena (default 60)

Opciones varias:
  --whitelist <MAC,...> Lista de BSSIDs en blanco (separados por coma)
  --help                Muestra esta ayuda
\end{lstlisting}

\textbf{Ejemplos de uso:}

\begin{lstlisting}
# Captura en vivo
sudo python3 detector.py --iface panda0 --channel 6 \
    --bssid e4:ab:89:d6:9b:80 --ssid LAB_WARDEN_UTP

# Analisis offline de captura previamente generada
python3 detector.py --pcap captures/chain-complete-sample.pcapng \
    --bssid e4:ab:89:d6:9b:80 --ssid LAB_WARDEN_UTP

# Captura con umbrales personalizados
sudo python3 detector.py --iface panda0 --channel 6 \
    --bssid e4:ab:89:d6:9b:80 --ssid LAB_WARDEN_UTP \
    --umbral-beacons 50 --umbral-deauth 8
\end{lstlisting}

\subsection{Diseño del Servidor FastAPI del Panel del Defensor}

El Panel del Defensor se sirve mediante un servidor FastAPI que integra
el detector como componente interno. Esta sección especifica los
endpoints HTTP, el protocolo WebSocket y la estructura del frontend.

\subsubsection{Estructura del proyecto Python}

\begin{lstlisting}
detector/
+-- main.py                    (entry point CLI clasico)
+-- web/
|   +-- __init__.py
|   +-- server.py              (aplicacion FastAPI + lifecycle)
|   +-- websocket_manager.py   (gestor de conexiones WS)
|   +-- routes.py              (endpoints HTTP REST)
|   +-- detector_runner.py     (hilo del detector, integracion)
|   +-- static/
|       +-- index.html         (entry point del frontend)
|       +-- panel.js
|       +-- styles.css
+-- capture/
+-- analyzers/
+-- correlator.py
+-- reporter.py                (modificado: emite a stdout y a callbacks)
+-- session.py                 (modificado: soporta callbacks de alerta)
+-- config.py
+-- requirements.txt
\end{lstlisting}

\subsubsection{Aplicación FastAPI: server.py}

\begin{lstlisting}
from fastapi import FastAPI, WebSocket, WebSocketDisconnect
from fastapi.staticfiles import StaticFiles
from fastapi.responses import FileResponse
from .websocket_manager import WebSocketManager
from .detector_runner import DetectorRunner
from .routes import router as api_router

app = FastAPI(title="WARDEN Panel del Defensor")

# Estado compartido del servidor
ws_manager = WebSocketManager()
detector_runner = DetectorRunner(ws_manager)

# Servir frontend estatico
app.mount("/static", StaticFiles(directory="web/static"), name="static")

@app.get("/")
async def root():
    return FileResponse("web/static/index.html")

# Endpoint WebSocket para alertas en vivo
@app.websocket("/ws")
async def websocket_endpoint(ws: WebSocket):
    await ws_manager.connect(ws)
    try:
        while True:
            # El frontend envia comandos (reset, etc.)
            mensaje = await ws.receive_json()
            await procesar_comando_ws(mensaje)
    except WebSocketDisconnect:
        ws_manager.disconnect(ws)

# Endpoints REST
app.include_router(api_router, prefix="/api")

# Lifecycle: detener detector cuando el servidor cierra
@app.on_event("shutdown")
async def shutdown():
    detector_runner.stop()
\end{lstlisting}

\subsubsection{Endpoints REST: routes.py}

\begin{lstlisting}
from fastapi import APIRouter

router = APIRouter()

@router.get("/status")
async def get_status():
    return {
        "detector_corriendo": detector_runner.is_running(),
        "duracion_seg": detector_runner.duracion_actual(),
        "alertas_totales": detector_runner.contadores_totales(),
        "frames_procesados": detector_runner.frames_procesados()
    }

@router.post("/detector/start")
async def start_detector(config: dict):
    # config: { iface, channel, bssid_protegido, ssid_protegido, umbrales }
    try:
        detector_runner.start(config)
        return { "ok": True }
    except Exception as e:
        return { "ok": False, "error": str(e) }

@router.post("/detector/stop")
async def stop_detector():
    detector_runner.stop()
    return { "ok": True }

@router.post("/session/reset")
async def reset_session():
    detector_runner.reset_session()
    await ws_manager.broadcast({ "tipo": "session_reset" })
    return { "ok": True }
\end{lstlisting}

\subsubsection{Gestor de WebSockets: websocket_manager.py}

\begin{lstlisting}
class WebSocketManager:
    def __init__(self):
        self.conexiones: list[WebSocket] = []

    async def connect(self, ws: WebSocket):
        await ws.accept()
        self.conexiones.append(ws)
        # Enviar estado inicial al conectarse
        await ws.send_json({
            "tipo": "init",
            "alertas_recientes": detector_runner.alertas_recientes(),
            "estado": "conectado"
        })

    def disconnect(self, ws: WebSocket):
        if ws in self.conexiones:
            self.conexiones.remove(ws)

    async def broadcast(self, mensaje: dict):
        # Envia mensaje a todos los clientes conectados
        muertos = []
        for ws in self.conexiones:
            try:
                await ws.send_json(mensaje)
            except Exception:
                muertos.append(ws)
        for ws in muertos:
            self.disconnect(ws)
\end{lstlisting}

\subsubsection{Hilo del detector: detector_runner.py}

\begin{lstlisting}
import threading
import asyncio
from queue import Queue

class DetectorRunner:
    def __init__(self, ws_manager):
        self.ws_manager = ws_manager
        self.detector_thread = None
        self.cola_alertas = Queue()
        self.contadores = {
            "BEACON_FLOOD": 0, "DEAUTH": 0,
            "EVIL_TWIN": 0, "CADENA_OFENSIVA": 0
        }
        self.frames = 0
        self.tiempo_inicio = None

    def start(self, config):
        # Validar interfaz en modo monitor
        verificar_modo_monitor(config["iface"])

        # Crear sesion del detector configurada para emitir alertas
        # tanto a stdout (Reporter) como a la cola (callback)
        sesion = ManejadorSesion(config)
        sesion.set_alert_callback(self._on_alert)

        # Lanzar en hilo separado
        self.detector_thread = threading.Thread(
            target=sesion.iniciar,
            daemon=True
        )
        self.tiempo_inicio = time.time()
        self.detector_thread.start()

        # Iniciar tarea asyncio que consume la cola
        asyncio.create_task(self._consumir_cola())

    def _on_alert(self, alerta):
        # Llamado desde el hilo del detector
        self.cola_alertas.put(alerta)
        if alerta["tipo"] in self.contadores:
            self.contadores[alerta["tipo"]] += 1

    async def _consumir_cola(self):
        # Tarea asyncio que reenvia alertas al WebSocket
        while True:
            if not self.cola_alertas.empty():
                alerta = self.cola_alertas.get_nowait()
                await self.ws_manager.broadcast({
                    "tipo": "alerta",
                    "datos": alerta
                })
            await asyncio.sleep(0.05)  # 20 Hz polling de la cola

    def stop(self):
        if self.detector_thread and self.detector_thread.is_alive():
            # Mecanismo de detencion del detector
            sesion.detener()
            self.detector_thread.join(timeout=5)

    def reset_session(self):
        for k in self.contadores:
            self.contadores[k] = 0
        self.frames = 0
        self.tiempo_inicio = time.time()
\end{lstlisting}

\subsubsection{Modificación del Reporter para soportar callbacks}

El componente Reporter del detector se modifica ligeramente para emitir
alertas tanto a stdout (comportamiento original) como a un callback
opcional registrado por el ManejadorSesion:

\begin{lstlisting}
class Reporter:
    def __init__(self, sesion):
        self.sesion = sesion
        self.alert_callback = None  # nuevo: opcional

    def set_alert_callback(self, fn):
        self.alert_callback = fn

    def alerta(self, severidad, tipo, mensaje, detalles):
        timestamp = datetime.now().isoformat()
        alerta_obj = {
            "timestamp": timestamp,
            "severidad": severidad,
            "tipo": tipo,
            "mensaje": mensaje,
            "detalles": detalles
        }

        # Emitir a stdout (comportamiento original, RF-DEF-008)
        self._imprimir_consola(alerta_obj)

        # Notificar al sesion para correlador y contadores
        self.sesion.registrar_alerta_emitida(tipo)

        # Si hay callback (modo Panel), invocarlo
        if self.alert_callback:
            self.alert_callback(alerta_obj)
\end{lstlisting}

\subsubsection{Protocolo WebSocket: mensajes}

Los mensajes que viajan por el WebSocket entre frontend y backend siguen
este esquema:

\textbf{Backend hacia frontend:}

\begin{lstlisting}
// Mensaje de inicio (al conectarse)
{
  "tipo": "init",
  "alertas_recientes": [...],
  "estado": "conectado"
}

// Alerta nueva
{
  "tipo": "alerta",
  "datos": {
    "timestamp": "2026-04-25T14:23:01",
    "severidad": "ALERT",
    "tipo": "BEACON_FLOOD",
    "mensaje": "Tasa anomala de beacons: 87.4 BSSIDs unicos/s",
    "detalles": { "bssids_unicos": 437, "ventana_seg": 5 }
  }
}

// Reset de sesion (despues de POST /api/session/reset)
{
  "tipo": "session_reset"
}

// Cambio de estado del detector
{
  "tipo": "detector_status",
  "estado": "corriendo" | "detenido" | "error",
  "mensaje": "..."
}
\end{lstlisting}

\textbf{Frontend hacia backend:}

\begin{lstlisting}
// Solicitar estado actual
{ "comando": "status" }

// Mantener heartbeat (keep-alive)
{ "comando": "ping" }
\end{lstlisting}

\subsubsection{Frontend: estructura HTML}

El frontend del Panel del Defensor es una página estática servida por
FastAPI. Estructura mínima:

\begin{lstlisting}
<!DOCTYPE html>
<html lang="es">
<head>
  <meta charset="UTF-8">
  <title>WARDEN - Panel del Defensor</title>
  <script src="https://cdn.tailwindcss.com"></script>
</head>
<body class="bg-gray-100 min-h-screen">
  <header class="bg-gray-900 text-white p-4">
    <h1 class="text-2xl font-bold">WARDEN - Panel del Defensor</h1>
  </header>

  <main class="max-w-6xl mx-auto p-4 grid grid-cols-3 gap-4">

    <!-- Indicador visual de amenaza -->
    <section class="col-span-3 bg-white rounded shadow p-6 text-center">
      <div id="threat-indicator" class="w-32 h-32 rounded-full mx-auto bg-green-500"></div>
      <p id="threat-label" class="mt-2 text-xl font-semibold">SIN AMENAZAS</p>
    </section>

    <!-- Contadores agregados -->
    <section class="col-span-3 bg-white rounded shadow p-4 grid grid-cols-4 gap-2">
      <div class="text-center">
        <p class="text-3xl" id="cnt-beacon">0</p>
        <p class="text-sm text-gray-600">Beacon Flood</p>
      </div>
      <div class="text-center">
        <p class="text-3xl" id="cnt-deauth">0</p>
        <p class="text-sm text-gray-600">Deauthentication</p>
      </div>
      <div class="text-center">
        <p class="text-3xl" id="cnt-eviltwin">0</p>
        <p class="text-sm text-gray-600">Evil Twin</p>
      </div>
      <div class="text-center">
        <p class="text-3xl text-red-600" id="cnt-cadena">0</p>
        <p class="text-sm text-gray-600">Cadena Ofensiva</p>
      </div>
    </section>

    <!-- Lista de alertas -->
    <section class="col-span-3 bg-white rounded shadow p-4">
      <div class="flex justify-between items-center">
        <h2 class="text-lg font-bold">Alertas en tiempo real</h2>
        <button onclick="resetSession()" class="bg-gray-200 hover:bg-gray-300 px-4 py-1 rounded">
          Reiniciar sesion
        </button>
      </div>
      <ul id="alertas-list" class="mt-3 space-y-2 max-h-96 overflow-y-auto">
        <!-- Las alertas se insertan aqui dinamicamente -->
      </ul>
    </section>

  </main>

  <script src="/static/panel.js"></script>
</body>
</html>
\end{lstlisting}

\subsubsection{Frontend: lógica JavaScript (panel.js)}

\begin{lstlisting}
// Conexion WebSocket
const ws = new WebSocket("ws://" + location.host + "/ws");

ws.onopen = () => console.log("Conectado al backend");
ws.onclose = () => {
  console.log("Conexion perdida, reintentando...");
  setTimeout(() => location.reload(), 3000);
};

ws.onmessage = (evt) => {
  const msg = JSON.parse(evt.data);
  switch (msg.tipo) {
    case "init":
      cargarAlertasIniciales(msg.alertas_recientes);
      break;
    case "alerta":
      agregarAlerta(msg.datos);
      actualizarContador(msg.datos.tipo);
      actualizarIndicador();
      break;
    case "session_reset":
      limpiarVista();
      break;
  }
};

function agregarAlerta(alerta) {
  const li = document.createElement("li");
  const colorClass = severidadAColor(alerta.severidad);
  li.className = "p-3 rounded border-l-4 " + colorClass;
  li.innerHTML = `
    <div class="flex justify-between">
      <span class="font-semibold">${alerta.tipo}</span>
      <span class="text-sm text-gray-500">${alerta.timestamp}</span>
    </div>
    <p class="mt-1">${alerta.mensaje}</p>
  `;
  // Insertar al tope (mas reciente primero)
  document.getElementById("alertas-list").prepend(li);
}

function actualizarIndicador() {
  const cnt = window.contadores;
  const indicator = document.getElementById("threat-indicator");
  const label = document.getElementById("threat-label");

  if (cnt.CADENA_OFENSIVA > 0) {
    indicator.className = "w-32 h-32 rounded-full mx-auto bg-red-600 animate-pulse";
    label.textContent = "CADENA OFENSIVA DETECTADA";
  } else if (cnt.BEACON_FLOOD + cnt.DEAUTH + cnt.EVIL_TWIN > 0) {
    indicator.className = "w-32 h-32 rounded-full mx-auto bg-yellow-500";
    label.textContent = "AMENAZA DETECTADA";
  } else {
    indicator.className = "w-32 h-32 rounded-full mx-auto bg-green-500";
    label.textContent = "SIN AMENAZAS";
  }
}

async function resetSession() {
  if (!confirm("Limpiar todas las alertas? El detector seguira corriendo.")) return;
  await fetch("/api/session/reset", { method: "POST" });
}
\end{lstlisting}

% ==============================================================================
\section{Diagramas de Secuencia de Flujos Críticos}
% ==============================================================================

\subsection{Flujo: Inicio de cadena automática}

\vspace{0.3cm}
\begin{verbatim}
Operador      Interfaz      Controlador   Modulo       Modulo      Modulo
Atacante      Serial        Cadena        Beacon       Deauth      EvilTwin
   |             |              |             |            |            |
   | ``start''     |              |             |            |            |
   |------------>|              |             |            |            |
   |             | iniciar()    |             |            |            |
   |             |------------->|             |            |            |
   |             |              |             |            |            |
   |             |              | iniciar_BF()|            |            |
   |             |              |------------>|            |            |
   |             |              |             |            |            |
   |             |              |  (30 seg)   |            |            |
   |             |              |             |            |            |
   |             |              |    fin_BF   |            |            |
   |             |              |<------------|            |            |
   |             |              |             |            |            |
   |             |              | iniciar_DA()|            |            |
   |             |              |--------------------------->|          |
   |             |              |             |            |            |
   |             |              |             |  (15 seg)  |            |
   |             |              |             |            |            |
   |             |              |    fin_DA   |            |            |
   |             |              |<---------------------------|          |
   |             |              |             |            |            |
   |             |              | iniciar_ET()|            |            |
   |             |              |---------------------------------------->|
   |             |              |             |            |            |
   |             |              |             |            | (120 seg)  |
   |             |              |             |            |            |
   |             |              |    fin_ET   |            |            |
   |             |              |<----------------------------------------|
   |             |              |             |            |            |
   |             |   resumen    |             |            |            |
   |             |<-------------|             |            |            |
   |  resumen    |              |             |            |            |
   |<------------|              |             |            |            |
\end{verbatim}

\subsection{Flujo: Procesamiento de un frame en el detector}

\vspace{0.3cm}
\begin{verbatim}
Scapy        Capturador    Despachador   Analizadores   Correlador   Reporter
  |              |              |              |              |          |
  | frame        |              |              |              |          |
  |------------->|              |              |              |          |
  |              | callback(f)  |              |              |          |
  |              |------------->|              |              |          |
  |              |              | filtro_canal |              |          |
  |              |              | tipo_frame   |              |          |
  |              |              |              |              |          |
  |              |              |  procesar(f) |              |          |
  |              |              |------------->|              |          |
  |              |              |              | analiza      |          |
  |              |              |              | ventanas     |          |
  |              |              |              |              |          |
  |              |              |              | si firma:    |          |
  |              |              |              |   alerta(..)>|          |
  |              |              |              |              | imprimir |
  |              |              |              |              |--------->|
  |              |              |              |              |          |
  |              |              |              |              | reg_alert|
  |              |              |              |              |<---------|
  |              |              |              |              | evaluar  |
  |              |              |              |              | cadena   |
  |              |              |              |              |          |
  |              |              |              |              | si 3 fases|
  |              |              |              |              | en orden:|
  |              |              |              |              | alerta()>|
  |              |              |              |              |          |
\end{verbatim}

\subsection{Flujo: Captura de credenciales en el portal cautivo}

\vspace{0.3cm}
\begin{verbatim}
Victima      Modulo        Servidor     Portal       Buffer
(navegador)  EvilTwin      DNS          HTTP         Credenciales
   |             |             |            |             |
   | DNS query   |             |            |             |
   | google.com  |             |            |             |
   |------------>|             |            |             |
   |             | resolver()  |            |             |
   |             |------------>|            |             |
   |             | 10.0.0.1    |            |             |
   |             |<------------|            |             |
   | DNS reply   |             |            |             |
   |<------------|             |            |             |
   |                                                       |
   | GET / HTTP/1.1                                        |
   | (a 10.0.0.1)                                          |
   |---------------------------------------->|             |
   |                                          | servir HTML |
   |                                          | portal.html |
   |<----------------------------------------|             |
   | (renderiza pagina)                                    |
   | usuario llena formulario                              |
   |                                                       |
   | POST /captura                                         |
   | usuario=demo.user@warden.test                         |
   | password=demo-password-12345                          |
   |---------------------------------------->|             |
   |                                          | parsear     |
   |                                          | construir   |
   |                                          | CredCapt    |
   |                                          |------------>|
   |                                          |             | almacenar
   |                                          |             | en buffer
   |                                          | log serial  |
   |                                          | imprimir    |
   |                                          | en consola  |
   |                                          |             |
   |                                          | servir HTML |
   |                                          | de exito    |
   |<----------------------------------------|             |
\end{verbatim}

\subsection{Flujo: Manejo de SIGINT en el detector}

\vspace{0.3cm}
\begin{verbatim}
Operador      Sistema      Manejador    Capturador    Reporter    Consola
Detector      Operativo    Sesion
   |             |             |             |             |          |
   | Ctrl+C      |             |             |             |          |
   |------------>|             |             |             |          |
   |             | SIGINT      |             |             |          |
   |             |------------>|             |             |          |
   |             |             | shutdown()  |             |          |
   |             |             |             |             |          |
   |             |             | detener()   |             |          |
   |             |             |------------>|             |          |
   |             |             |     ok      |             |          |
   |             |             |<------------|             |          |
   |             |             |             |             |          |
   |             |             | resumen()   |             |          |
   |             |             |---------------------------> |        |
   |             |             |             |             | imprimir |
   |             |             |             |             |--------->|
   |             |             |             |             |          |
   |             |             | sys.exit(0) |             |          |
   |             |<------------|             |             |          |
   | <prompt>    |             |             |             |          |
   |<------------|             |             |             |          |
\end{verbatim}

\subsection{Flujo: Lanzamiento de ataque desde el Panel del Atacante}

\vspace{0.3cm}
\begin{verbatim}
Operador     Panel        ESP32         Validador    Controlador  Modulos
Atacante     (browser)    (API REST)    Etico        Cadena       de ataque
   |             |              |             |             |          |
   |  click      |              |             |             |          |
   | "Lanzar"    |              |             |             |          |
   |------------>|              |             |             |          |
   |             | POST /config |             |             |          |
   |             | {bssid,mac}  |             |             |          |
   |             |------------->|             |             |          |
   |             |              | validar()   |             |          |
   |             |              |------------>|             |          |
   |             |              |    OK       |             |          |
   |             |              |<------------|             |          |
   |             |   200 OK     |             |             |          |
   |             |<-------------|             |             |          |
   |             |              |             |             |          |
   |             | POST /attack/start         |             |          |
   |             | {modo:cadena}              |             |          |
   |             |------------->|             |             |          |
   |             |              | iniciar()   |             |          |
   |             |              |---------------------------> |        |
   |             |              |             |             |          |
   |             |              |             |             | iniciar  |
   |             |              |             |             | beacon() |
   |             |              |             |             |--------->|
   |             |   200 OK     |             |             |          |
   |             |<-------------|             |             |          |
   |             |              |             |             |          |
   |             | (transita a vista de monitoreo)          |          |
   |             |              |             |             |          |
   |             | GET /status  |             |             |          |
   |             | (cada 1s)    |             |             |          |
   |             |------------->|             |             |          |
   |             |   estado JSON|             |             |          |
   |             |<-------------|             |             |          |
   |             |              |             |             |          |
   |             | (refresca contadores en pantalla)        |          |
   |             |              |             |             |          |
\end{verbatim}

\subsection{Flujo: Alerta del detector hacia el Panel del Defensor}

\vspace{0.3cm}
\begin{verbatim}
Capturador   Despachador  Analizador   Reporter   WebSocket   Frontend
Frames                    (X)                     Manager    (browser)
   |              |             |            |          |          |
   | frame        |             |            |          |          |
   |------------->|             |            |          |          |
   |              | enrutar     |            |          |          |
   |              |------------>|            |          |          |
   |              |             | si firma:  |          |          |
   |              |             | construir  |          |          |
   |              |             | alerta     |          |          |
   |              |             |            |          |          |
   |              |             | alerta(..) |          |          |
   |              |             |----------->|          |          |
   |              |             |            | imprimir |          |
   |              |             |            | a stdout |          |
   |              |             |            |          |          |
   |              |             |            | callback |          |
   |              |             |            | (cola)   |          |
   |              |             |            |--------->|          |
   |              |             |            |          | broadcast|
   |              |             |            |          | a todos  |
   |              |             |            |          | clientes |
   |              |             |            |          |--------->|
   |              |             |            |          |          | render
   |              |             |            |          |          | alerta
   |              |             |            |          |          | en lista
   |              |             |            |          |          |
   |              |             |            |          |          | actualizar
   |              |             |            |          |          | indicador
   |              |             |            |          |          | (verde->
   |              |             |            |          |          |  amarillo)
\end{verbatim}

% ==============================================================================
\section{Consideraciones de Implementación}
% ==============================================================================

\subsection{Consideraciones generales}

\begin{itemize}
  \item \textbf{Manejo de errores:} cada operación que pueda fallar
    (apertura de interfaz, lectura de archivo, parseo de frame) debe
    estar protegida con manejo de excepciones. Errores no críticos se
    logan y la ejecución continúa; errores críticos terminan
    ordenadamente con código de salida no cero.
  \item \textbf{Logging:} todos los componentes utilizan el módulo
    \texttt{logging} de Python con niveles apropiados (DEBUG, INFO,
    WARNING, ERROR). El nivel por defecto es INFO; se puede ajustar por
    argumento.
  \item \textbf{Tipos:} en Python se utilizan type hints en todas las
    firmas de funciones para facilitar la lectura y permitir
    verificación con \texttt{mypy} si se requiere.
\end{itemize}

\subsection{Subsistema ofensivo (ESP32)}

\begin{itemize}
  \item \textbf{Concurrencia:} el firmware utiliza FreeRTOS por debajo de
    Arduino-ESP32. Las tareas se crean con \texttt{xTaskCreate()} con
    prioridades cuidadosamente asignadas para no bloquear el WiFi.
  \item \textbf{Watchdog:} el watchdog del ESP32 debe ser servido
    regularmente en cualquier loop largo. Considerar
    \texttt{esp\_task\_wdt\_reset()} en bucles que excedan 1 segundo sin
    \texttt{delay()}.
  \item \textbf{Memoria:} la ESP32 tiene aproximadamente 320 KB de RAM
    útil. El buffer de credenciales y el HTML embebido del portal cabin
    sin problema.
  \item \textbf{Inyección de frames:} la inyección de frames 802.11 de
    gestión arbitrarios no es soportada por el SDK oficial de Espressif.
    Se requiere usar \texttt{esp\_wifi\_80211\_tx()} (función no
    documentada pero disponible) o bibliotecas de la comunidad como
    \texttt{esp32-wifi-penetration-tool}.
\end{itemize}

\subsection{Subsistema defensivo (Python)}

\begin{itemize}
  \item \textbf{Permisos:} la captura en vivo requiere
    \texttt{CAP\_NET\_RAW} y \texttt{CAP\_NET\_ADMIN}. Se ejecuta con
    \texttt{sudo} o se otorgan capabilities al binario de Python con
    \texttt{setcap}.
  \item \textbf{Rendimiento:} Scapy puede tener cuellos de botella ante
    altas tasas de frames. Para condiciones extremas (más de 1000
    frames/segundo), considerar el uso de \texttt{tcpdump} para captura
    a archivo y procesamiento offline en bloques.
  \item \textbf{Compatibilidad de Scapy con 802.11:} verificar que el
    formato de captura incluye el encabezado RadioTap. Algunas
    distribuciones requieren configuración adicional del adaptador para
    incluirlo.
  \item \textbf{Falsos positivos en redes complejas:} en entornos con
    múltiples APs legítimos compartiendo SSID (redes empresariales con
    roaming), la lista blanca de BSSIDs (RF-DEF-015) es esencial.
\end{itemize}

% ==============================================================================
\section{Trazabilidad: Diseño a Casos de Uso}
% ==============================================================================

\begin{datatable}{L{4cm} L{11.3cm}}{Componente de Diseño & Casos de Uso Implementados}
\drow{Controlador de Cadena (off)   & UC-01, UC-05, UC-16, UC-17}
\drow{Módulo Beacon Flood (off)     & UC-02 (y UC-01 fase 1, UC-16 fase 1)}
\drow{Módulo Deauth (off)           & UC-03 (y UC-01 fase 2, UC-16 fase 2)}
\drow{Módulo Evil Twin (off)        & UC-04 (y UC-01 fase 3, UC-16 fase 3)}
\drow{Portal Cautivo (off)          & UC-04, UC-11, UC-17 (visualización de credenciales)}
\drow{Validador Ético (off)         & UC-06, UC-15, UC-16}
\drow{Interfaz Serial (off)         & UC-01 a UC-06 (mecanismo secundario de control)}
\drow{AP de Control (off)           & UC-13 a UC-17 (puerta de entrada al panel)}
\drow{Servidor API REST (off)       & UC-13 a UC-17}
\drow{Reconocedor (off)             & UC-13 (\texttt{/scan}), UC-14 (\texttt{/clients})}
\drow{Lookup OUI (off)              & UC-14 (\texttt{/oui-lookup})}
\drow{Frontend Panel del Atacante   & UC-13 a UC-17}
\drow{Capturador de Frames (def)    & UC-07, UC-08}
\drow{Despachador (def)             & UC-07, UC-08}
\drow{Analizador Beacon Flood (def) & UC-07, UC-08, UC-09, UC-19}
\drow{Analizador Deauth (def)       & UC-07, UC-08, UC-09, UC-19}
\drow{Analizador Evil Twin (def)    & UC-07, UC-08, UC-09, UC-19}
\drow{Correlador de Cadena (def)    & UC-09, UC-19 (alerta agregada)}
\drow{Reporter (def)                & UC-09, UC-19 (callback hacia WebSocket)}
\drow{Manejador de Sesión (def)     & UC-07, UC-08, UC-10}
\drow{Servidor FastAPI (def)        & UC-18, UC-19, UC-20}
\drow{WebSocket Manager (def)       & UC-19 (alertas tiempo real), UC-20 (broadcast reset)}
\drow{Detector Runner (def)         & UC-18, UC-20}
\drow{Frontend Panel del Defensor   & UC-18, UC-19, UC-20}
\drow{Demostración integrada        & UC-12}
\end{datatable}

\end{document}
